%%%%%%%%%%%%%%%%%%%%%%%%%%%%%%%%%%%%%%%%%%%%%%%%%%%%%%%%%%%%%%%%%%%%%%%
% Universidade Federal de Santa Catarina
% Biblioteca Universitária
%----------------------------------------------------------------------
% Exemplo de utilização da documentclass ufscThesis
%----------------------------------------------------------------------
% (c)2013 Roberto Simoni (roberto.emc@gmail.com)
%         Carlos R Rocha (cticarlo@gmail.com)
%         Rafael M Casali (rafaelmcasali@yahoo.com.br)
%%%%%%%%%%%%%%%%%%%%%%%%%%%%%%%%%%%%%%%%%%%%%%%%%%%%%%%%%%%%%%%%%%%%%%%
\documentclass{ufscThesis} % Definicao do documentclass ufscThesis	

%----------------------------------------------------------------------
% Pacotes usados especificamente neste documento
\usepackage{graphicx} % Possibilita o uso de figuras e gráficos
\usepackage{color}    % Possibilita o uso de cores no documento
\usepackage{listings}
% The package listings is unable to deal correctly
% with multibyte strings.
% A workaround is to replace every "offending character"
% with its respective TeX escape codes.
%
% This is acomplished with the 'literate programming'
% feature of the listings package.
\lstset{
literate=
    {á}{{\' a}}1
    {Á}{{\' A}}1
    {à}{{\` a}}1
    {À}{{\` A}}1
    {â}{{\^ a}}1
    {Â}{{\^ A}}1
    {ã}{{\~ a}}1
    {Ã}{{\~ A}}1
    {ç}{{\c c}}1
    {Ç}{{\c C}}1
    {é}{{\' e}}1
    {É}{{\' E}}1
    {ê}{{\^ e}}1
    {Ê}{{\^ E}}1
    {í}{{\' \i}}1
    {Í}{{\' I}}1
    {ó}{{\' o}}1
    {Ó}{{\' O}}1
    {ô}{{\^ o}}1
    {Ô}{{\^ O}}1
    {õ}{{\~ o}}1
    {Õ}{{\~ O}}1
    {ú}{{\' u}}1
    {Ú}{{\' U}}1
}

%----------------------------------------------------------------------
% Comandos criados pelo usuário
\newcommand{\afazer}[1]{{\color{red}{#1}}} % Para destacar uma parte a ser trabalhada

%----------------------------------------------------------------------
% Início do documento
\begin{document}

%--------------------------------------------------------
% Elementos pré-textuais
% Capa e folha de rosto

\instituicao[a]{Universidade Federal de Santa Catarina} % Opcional
\departamento[a]{Biblioteca Universitária}
\curso[o]{Programa de ...}
\documento[a]{Tese} % [o] para dissertação [a] para tese
\titulo{Título}
\subtitulo{Subtítulo (se houver)} % Opcional
\autor{Nome completo do autor}
\grau{...}
\local{Florianópolis} % Opcional (Florianópolis é o padrão)
\data{04}{junho}{2013}
\orientador[Orientador\\Universidade ...]{Prof. Dr.}
\coorientador[Coorientador\\Universidade ...]{Prof. Dr.}
\coordenador[Coordenador\\Universidade ...]{Prof. Dr. }

\capa
\folhaderosto[comficha] % Se nao quiser imprimir a ficha, é só não usar o parâmetro

\numerodemembrosnabanca{3}
\bancaMembroA{
    Profa. Dra. Karina Girardi Roggia \\
    Universidade do Estado de Santa Catarina
}
\bancaMembroB{
    Prof. Dr. Ricardo Azambuja Silveira \\
    Universidade Federal de Santa Catarina
}
\bancaMembroC{
    Prof. Dr. Rosvelter João Coelho da Costa \\
    Universidade Federal de Santa Catarina
}

\folhaaprovacao

% Dedicatória

\dedicatoria{
    À minha mãe, Venida
}

\paginadedicatoria

% Agradecimentos

\agradecimento{
    Agradeço muito minha família, em especial à minha mãe,
    por sempre terem me apoiado, guiado e auxiliado
    ao longo dos últimos 20 anos.

    Agradeço minha orientadora, Jerusa,
    que me acolheu logo no início do curso,
    por todas as conversas e orientações,
    que muitas vezes extrapolaram os limites da Academia.

    Agradeço aos professores do Programa Avançado de Matemática
    (Melissa, Gilles, Fernando),
    que me ajudaram a construir a maturidade matemática
    que me foi extremamente útil nos quatro anos de graduação.

    Agradeço as discussões com Santana, Schultz, Gabriel e Abouhatem
    sobre os algoritmos da Maratona de Programação.

    Agradeço todas as discussões produtivas (e improdutivas)
    com o pessoal do laboratório IATE.

    E, por último, agradeço o companheirismo dos meus amigos
    ao longo desta jornada.
}

\paginaagradecimento

% Epígrafe

\epigrafe{
    Even if one works basically all week trying to prove $\P \neq \NP$,
    one should set aside Friday afternoon for trying to prove $\P = \NP$.
}
{
    (John Edward Hopcroft)
}
\paginaepigrafe

% Resumo

\textoResumo{
    O texto do resumo deve ser digitado,
    em um único bloco,
    sem espaço de parágrafo.
    O resumo deve ser significativo,
    composto de uma sequência de frases concisas,
    afirmativas e não de uma enumeração de tópicos.
    Não deve conter citações.
    Deve usar o verbo na voz passiva.
    Abaixo do resumo,
    deve-se informar as palavras-chave
    (palavras ou expressões significativas retiradas do texto)
    ou, termos retirados de thesaurus da área.
}
\palavrasChave{
    Palavra-chave 1.
    Palavra-chave 2.
    Palavra-chave 3.
}

\paginaresumo

% Resumo, em inglês.

\textAbstract{
    Resumo traduzido para outros idiomas,
    neste caso,
    inglês.
    Segue o formato do resumo
    feito na língua vernácula.
    As palavras-chave traduzidas,
    versão em língua estrangeira,
    são colocadas abaixo do texto
    precedidas pela expressão ``Keywords'',
    separadas por ponto.
}
\keywords{
    Keyword 1.
    Keyword 2.
    Keyword 3.
}

\paginaabstract

% Sumário e listas usadas no documento.
% As listas dependem da necessidade do usuário

\listadefiguras
\listadetabelas
\listadeabreviaturas
\listadesimbolos
\sumario


%--------------------------------------------------------
% Elementos textuais
\emph{Complexidade} é a quantidade de recursos
que uma máquina de Turing gasta
para computar determinada função
ou para decidir pertinência a uma linguagem
\cite[p.~285]{HopcroftUllman1979}.
Os recursos mais importantes para a teoria de complexidade computacional
são o espaço e o tempo,
em máquinas de Turing determinísticas e não"=determinísticas.

Os axiomas de Blum definem a noção de complexidade computacional
não apenas para máquinas de Turing,
mas sim para qualquer modelo de computação que satisfaça algumas restrições.

A seção~\ref{sec:enumeration_of_recursive_functions}
contém a maquinaria matemática necessária para expressar estas restrições.
Os axiomas de Blum aparecem na seção~\ref{sec:blum_axioms}.
Por fim, a seção~\ref{sec:default_measures}
define as medidas de complexidade padrão.

% Desenvolvimento
\chapter{EXPOSIÇÃO DO TEMA OU MATÉRIA}

É a parte principal e mais extensa do trabalho.
Deve apresentar a fundamentação teórica,
a metodologia,
os resultados
e a discussão.
Divide-se em seções e subseções
conforme a NBR 6024~\cite{abnt14724}.
Quanto a sua estrutura,
segue as recomendações da norma
para preparação de trabalhos acadêmicos,
a NBR 14724 de 2011~\cite{abnt14724}.
Quanto à Formatação,
segue o modelo adotado pela UFSC, o formato A5.

\section{FORMATAÇÃO DO TEXTO}

No que diz respeito à estrutura do trabalho,
o novo modelo para dissertações e teses
adotado pela UFSC
segue a NBR 14724 (2011).
Porém, em relação à formatação,
a UFSC adotou o tamanho A5,
que corresponde à metade do A4.
Por esta razão,
foi necessário uma adequação no tamanho da fonte,
espaçamento entrelinhas,
margens,
etc,
conforme exposto no quadro abaixo.

O texto deve ser justificado,
digitado em cor preta,
podendo utilizar outras cores
somente para as ilustrações.
Utilizar papel branco.
Os elementos pré-textuais
devem iniciar no anverso da folha,
com exceção da ficha catalográfica.
Os elementos textuais e pós-textuais
devem ser digitados no anverso e verso das folhas,
com espaçamento simples (1).

\subsection{As ilustrações}

Independente do tipo de ilustração 
(quadro, desenho, figura, fotografia, mapa, entre outros)
sua identificação aparece na parte superior,
precedida da palavra designativa.

A indicação da fonte consultada
deve aparecer na parte inferior,
elemento obrigatório mesmo que seja produção do próprio autor.
A ilustração deve ser citada no texto
e inserida o mais próximo possível do texto
a que se refere~\cite{abnt14724}.

A Figura \ref{fig:a} mostra um quadrado.
\begin{figure}[!htb]
   \centering
   \caption{Quadrado.}\label{fig:a}
   \begin{tikzpicture}
        \draw (0, 0) rectangle (1, 1);
    \end{tikzpicture}
\end{figure}

A Tabela~\ref{tab:a} mostra mais informações do template BU.

\begin{table}[!htb]
\begin{center}
    \caption{Formatação do texto.}\label{tab:a}
    \begin{tabular}{ p{3cm} | p{6cm} }
        \hline
        Cor & Branco \\
        \hline
        Formato do papel & A5 \\
        \hline
        Gramatura & 75 \\
        \hline
        Impressão & Frente e verso \\
        \hline
        Margens & Espelhadas: superior 2, Inferior: 1,5, Externa 1,5 e Externa: 2. \\
        \hline
        Cabeçalho & 0,7 \\
        \hline
        Rodapé & 0,7 \\
        \hline
        Paginação & Externa \\
        \hline
        Alinhamento vertical & Superior \\
        \hline
        Alinhamento do texto & Justificado \\
        \hline
        Fonte sugerida & Times New Roman  \\
        \hline
        Tamanho da fonte & 10,5 para o texto
        incluindo os títulos das seções e subseções.
        As citações com mais de três linhas
        as legendas das ilustrações e tabelas, fonte 9,5. \\
        \hline
        Espaçamento entre linhas & Um (1) simples \\
        \hline
        Espaçamento entre parágrafos & Anterior 0,0; Posterior 0,0 \\
        \hline
        Numeração da seção & As seções primárias devem começar
        sempre em páginas ímpares.
        Deixar um espaço (simples) entre o título da seção
        e o texto e entre o texto e o título da subseção. \\
        \hline
    \end{tabular}
\end{center}
Fonte: Universidade Federal de Santa Catarina (2011)
\end{table}

\subsection{Equações e fórmulas}

As equações e fórmulas
devem ser destacadas no texto
para facilitar a leitura.
Para numerá-las,
deve-se usar algarismos arábicos entre parênteses
e alinhados à direita.
Pode-se usar uma entrelinha maior do que a usada no texto~\cite{abnt14724}.

Exemplo: A equação \ref{eq:a}
\begin{equation}
    x^2 + y^2 = z^2
    \label{eq:a}
\end{equation}
 e a equação \ref{eq:b}
\begin{equation}
    x^2 + y^2 = n
    \label{eq:b}
\end{equation}

\section{EXEMPLO DE COMO GERAR A LISTA DE SÍMBOLOS E ABREVIATURAS}

Para gerar a lista de símbolos e abreviaturas use os comandos

\abreviatura{ABNT}{Associação Brasileira de Normas Técnicas}
\abreviatura{IBGE}{Instituto Brasileiro de Geografia e Estatística}
\simbolo{$\int$}{Integral}
\simbolo{$\prod$}{Produtório}

\begin{lstlisting}
\simbolo{símbolo}{descrição}
\end{lstlisting}

\begin{lstlisting}
\abreviatura{abreviatura}{descrição}
\end{lstlisting}

\subsection{Exemplo de citações no \LaTeX}

Segundo \citeonline{alves_2001} ...

...no final da frase \cite{abnt14724,BU_formatoA5}

\nocite{alves_2001,abnt10520,abnt6024,abnt14724}
 % Aqui vêm uma sequência de capítulos
\chapter{Comparação com outros trabalhos}

Este capítulo discute brevemente outros trabalhos que,
de certa forma,
tentam capturar a noção de ``função não"=determinística''.

\section{Problemas de busca vs problemas de decisão (Papadimitriou)}
\label{sec:papadimitriou_comparison}

\citeonline[p.~229]{Papadimitriou1994}
define $\FNP$ diferente de nossa definição
(na seção~\ref{sec:functional_complexity}).

\begin{definition}
    Uma linguagem $L$ é \emph{polinomialmente equilibrada}
    se existir algum polinômio $p$ tal que
    todos os elementos de $L$ são pares ordenados da forma $(x, y)$
    em que $|y| \leq p(|x|)$.
\end{definition}
Isto é, $L$ é constituída de pares,
de forma que o segundo elemento do par não é muito maior que o primeiro elemento.

Linguagens de $\NP$ e linguagens polinomialmente equilibradas
estão relacionadas por certificados de pertinência.
Por exemplo,
para a linguagem $\SAT$,
podemos provar eficientemente
(isto é, em tempo polinomial)
que determinada instância $\varphi$ é satisfazível
se fornecermos uma atribuição de valores"=verdade que satisfaz a instância;
esta atribuição é uma \emph{testemunha} ou \emph{certificado}
para a pertinência de $\varphi$ a $\SAT$.

Para cada linguagem $L \in \NP$,
podemos sistematicamente prover certificados de pertinência para $L$:
como sabemos que existe uma máquina de Turing não"=determinística que decide $L$,
podemos fornecer a sequência de transições não"=determinísticas
como um certificado de pertinência.
Como esta máquina opera em tempo polinomial,
para cada $x \in L$,
o par
\begin{equation*}
    (x, y),
\end{equation*}
em que $y$ é esta sequência de transições,
satisfará $|y| \leq p(|x|)$,
para algum polinômio $p$
--- no caso, $p$ é o próprio limite de tempo da máquina não"=determinística
que reconhece $L$.

Dessa forma,
podemos, sistematicamente,
associar uma linguagem $L \in \NP$
a uma linguagem polinomialmente equilibrada $R_L \in \P$.
Assim,
construiremos o conjunto $\FNP$ definido por \citeonline[p.~229]{Papadimitriou1994}.

\begin{definition}
    Se $L \in \NP$, chame de $R_L$
    uma linguagem polinomialmente equilibrada associada com $L$.
    Então defina $\FNP$ por
    \begin{equation*}
        \FNP = \{R_L \mid L \in \NP\}
    \end{equation*}
    \cite[p.~229]{Papadimitriou1994}
\end{definition}

Estes problemas são por vezes chamados de \emph{problemas de busca}
(do inglês \emph{search problem}),
em oposição a \emph{problemas de decisão}.

Do ponto de vista computacional,
a definição de Papadimitriou captura a ideia de
encontrar alguma solução para o problema $L$,
potencialmente descartando uma quantidade exponencial de outras soluções.

Esta definição contorna o problema de definir ``função não"=determinística''
trabalhando diretamente com os certificados de pertinência.
Entretanto,
a ausência do conceito de função
dificulta a formalização do conceito de oráculo
(pois não há mais uma única resposta ``certa'',
mas sim várias)
e impossibilita a análise do efeito da composição de funções,
como a feita na seção~\ref{sec:function_composition}.

\section{Conjunto potência como contradomínio (Andreev)}

\citeonline[p.~3]{Andreev1994}
apresenta uma abordagem diferente.
Em vez de tentar amarrar algum dos ramos de computação,
a função retornará todos eles.

\begin{definition}
    Sejam $\mathcal A$ e $\mathcal B$ conjuntos finitos,
    e $2^\mathcal B$ o conjunto potência de $\mathcal B$.
    Uma \emph{função não"=determinística de $\mathcal A$ em $\mathcal B$}
    é uma função da forma
    \begin{equation*}
        f : \mathcal A \to 2^\mathcal B.
    \end{equation*}
    \cite[p.~3]{Andreev1994}
\end{definition}

Os conjuntos $\mathcal A$ e $\mathcal B$ usados por \citeonline[p.~4]{Andreev1994}
são da forma $\{0, 1\}^k$.
A ideia é que $\mathcal A$ represente as possíveis entradas para um circuito booleano,
e que $\mathcal B$ represente as saídas desse circuito.
Dado um circuito computacional $S$,
com $k$ entradas e $l$ saídas,
dizemos que $S$ computa a função não"=determinística $F$ de $\{0, 1\}^k$ em $\{0, 1\}^l$
se,
para todo $a \in \{0, 1\}^k$,
$S(a) \in F(a)$.
Isto é,
todas as saídas possíveis de $S$ estão previstas em $F$.
Esta definição é similar à definição de \citeonline[p.~229]{Papadimitriou1994},
no sentido de que basta retornar um valor de $F(a)$.
\citeonline[p.~4]{Andreev1994}
define, então,
a \emph{complexidade} de uma função não"=determinística $F$
como sendo o tamanho do menor circuito que computa $F$.

O fato de Andreev se restringir a domínio e contradomínio finitos
mostra que esta definição não ``combina'' corretamente
a noção de função determinística com máquinas não"=determinísticas,
pois as funções operam sobre o domínio $\{0, 1\}^*$.

Além disso,
como o domínio e o contradomínio das funções não"=determinísticas de Andreev
são diferentes,
seria necessário um malabarismo adicional
para desenvolver a noção de composição de funções,
assim como na definição de Papadimitriou.

\section{Definição de Hopcroft e Ullman}

A definição de \citeonline[p.~313]{HopcroftUllman1979},
mencionada na introdução,
é a que motivou este trabalho.
\begin{definition}
    Seja $M$ uma máquina de Turing não"=determinística.
    Dizemos que $M(x) = y$ se
    algum dos ramos de computação de $M$ em $x$ produz $y$,
    e nenhum outro ramo de computação que se encerre
    produz um valor diferente de $y$.
    \cite[p.~313]{HopcroftUllman1979}
\end{definition}

O problema é que esta definição não enumera corretamente as funções recursivas.
O grande vilão está na permissão que damos à máquina $M$
não parar em todos os ramos de computação.

Por exemplo,
construa uma máquina $N$ que,
na entrada $\langle M, x \rangle$,
assuma dois ramos de computação:
no primeiro,
$N$ incondicionalmente escreve $1$ na fita de saída e para.
No outro, $N$ simula a máquina determinística $M$ em $x$ e,
caso $M$ pare ao computar $x$,
$N$ escreve $0$ na fita.

Se $M$ não para em $x$,
então $N$ possui um único ramo de computação que para,
e este ramo produz $1$ na fita;
portanto, $N(\langle M, x \rangle) = 1$.
Caso $M$ pare ao computar $x$,
$N(\langle M, x \rangle)$ estará indefinido,
pois dois ramos de computação de $N$ escrevem coisas diferentes na fita.

Ou seja,
esta máquina computa
\begin{equation*}
    f(\langle M, x \rangle) = 1, \text{ se $M$ não parar ao computar $x$,}
\end{equation*}
que é exatamente o complemento do problema da parada.
Portanto,
pela definição de Hopcroft e Ullman,
conseguimos enumerar funções não"=recursivas.

Mas,
se impormos à definição de Hopcroft e Ullman a restrição adicional
de que todos os ramos de $M$ devem parar para que $M(x)$ exista,
então não há a necessidade de haver vários ramos de computação,
pois todos eles retornam o mesmo resultado.
Basta fixar uma das possíveis transições não"=determinísticas,
tornando a máquina \emph{determinística}.
Neste caso,
perdemos o aparente ganho de tempo exponencial ao usar não"=determinismo.


\bibliographystyle{ufscThesis/ufsc-alf}
\bibliography{bibliografia}

%--------------------------------------------------------
% Elementos pós-textuais
\apendice
% Exemplo de apêndice

\chapter{EXEMPLIFICANDO UM APÊNDICE}

Texto do Apêndice aqui.


\anexo
% Exemplo de anexo.

\chapter{EXEMPLIFICANDO UM ANEXO}

Texto do anexo aqui.


\end{document}
